\section{Limitations}
\label{sec:limits}
% HR: reviewers will ask these anyway so better we say them first
\begin{itemize}
  \item Absolute numbers are low. Even with a perfect vocabulary the frozen \gdino-T reach only 0.107 AP50 on test, and some classes stay near zero.
  \item The gain is limited to imagery and label spaces close to RemoteCLIP's training data.
  \item One detector (Swin-T) and one benchmark family. We did not test bigger detectors like \gdino-B or OWLv2. \TODO{if time / compute allows}
  \item The rule family and the target size of 8 were designed by us on dev, and the synonym experiment was motivated by our own dev failure analysis. Both were frozen before test but readers should know.
\end{itemize}

\section{Conclusion}
\label{sec:conclusion}
% TODO: rewrite, this is just placeholder text from the meeting
We showed that choosing a small per-image vocabulary from pixels, with an aerial CLIP model and a calibrated adaptive size, robustly improves a frozen open-vocabulary detector on held-out aerial images at almost no cost.
The benefit is real but bounded: it holds where the encoder knows the imagery, and it is capped by near-synonym classes across merged datasets and by the detector on small objects.
\TODO{future work: taxonomy-aware vocabularies, aerial-adapted detector, other OVD models}
